% Fig 5:fig17 -- Estructura del prompt de extraccion (DeepSeek cache hit)
\begin{tikzpicture}[
    node distance=0.12cm,
    every node/.style={font=\scriptsize, align=left, inner sep=4pt},
    sys/.style={draw, rounded corners=2pt, fill=blue!10, minimum width=12cm, minimum height=0.7cm},
    inv/.style={draw, rounded corners=2pt, fill=green!12, minimum width=12cm, minimum height=0.55cm},
    var/.style={draw, rounded corners=2pt, fill=orange!12, minimum width=12cm, minimum height=0.55cm},
    note/.style={font=\scriptsize\itshape, align=left}
]
    \node[sys] (sys) {\textbf{SYSTEM} -- rol ``analista politico del Peru'' + contrato JSON estricto};
    \node[inv, below=of sys] (b1) {\textbf{USER (prefijo invariante -- cacheado)}};
    \node[inv, below=of b1] (b2) {Catalogo de candidatos (40)};
    \node[inv, below=of b2] (b3) {Categorias (17, las 5 ultimas no-politicas)};
    \node[inv, below=of b3] (b4) {Modos narrativos / aspectos / regiones / acciones};
    \node[inv, below=of b4] (b5) {Schema JSON de respuesta};
    \node[inv, below=of b5] (b6) {Ejemplos few-shot (5)};
    \node[inv, below=of b6] (b7) {Instrucciones (9 reglas, regla \#5 = no-politico)};

    \node[var, below=0.35cm of b7] (v1) {\textbf{USER (sufijo variable)}};
    \node[var, below=of v1] (v2) {Menu de temas del dia (5 tiers, hasta 20)};
    \node[var, below=of v2] (v3) {Documento a etiquetar (1 por llamada)};

    \node[note, right=0.2cm of b4] {prompt cache hit 85.5\%};
    \node[note, right=0.2cm of v2] {cambia por documento};
\end{tikzpicture}
